\documentclass[10pt,a4paper]{amsart}
\usepackage[margin=19mm]{geometry}
\usepackage{amsmath,amssymb,mathtools}
\usepackage[scr=rsfs,cal=euler]{mathalfa}
\usepackage{xfrac}
\usepackage[unicode,hidelinks,bookmarks=false]{hyperref}
\newcommand{\ie}{i.e.\ }
\newcommand{\cf}{cf.\ }
\newcommand{\resp}{respectively\ }
\newcommand{\Subsection}{\subsection}
\providecommand{\nobreakdash}{\nobreak\hbox{-}}
\setlength{\emergencystretch}{2em}
\multlinegap0pt
\usepackage{bm}
\usepackage{bbm}
\usepackage{dsfont}
\usepackage{xspace}
\usepackage{cancel}
\usepackage{mathtools}
\usepackage{amsthm}
\makeatletter
\@ifundefined{theorem}{\newtheorem{theorem}{Theorem}}{}
\@ifundefined{lemma}{\newtheorem{lemma}[theorem]{Lemma}}{}
\makeatother
\newcommand\calU{\mathcal{U}}
\newcommand\calZ{\mathcal{Z}}
\newcommand\clbeta[1][X]{\cl_{\beta #1}}
\newcommand\cl{\operatorname{cl}}
\newcommand\preim{^\gets}
\title[$C$-embedding, Lindel\"ofness, and \v{C}ech-completeness]{$C$-embedding, Lindel\"ofness, and \v{C}ech-completeness}
\author{Alan Dow, Klaas Pieter Hart, Jan van Mill, Hans Vermeer}
\date{Source: arXiv:2404.19703v3 (CC BY 4.0)}
\begin{document}
\maketitle

\noindent\emph{Source and licence.} $C$-embedding, Lindel\"ofness, and \v{C}ech-completeness. Version arXiv:2404.19703v3 (CC BY 4.0), distributed under the Creative Commons Attribution 4.0 International licence, as recorded in the arXiv licence metadata for this exact version and shown on the abstract page https://arxiv.org/abs/2404.19703v3. Full rights notice: licenses/arxiv-2404.19703-CC-BY-4.0.txt.\par\smallskip

\noindent\textit{Editorial scope.} Counted unit: the Theorem whose source label is \texttt{thm.first} in the article (source0.tex lines 272--284, source label \texttt{thm.first}), delivered together with the article's own complete original proof of it (source0.tex lines 286--352, one \texttt{proof} environment of 67 lines). No mathematics is written, repaired or re-derived: statement and proof are the article's own text line for line. Required context retained uncounted: lemma.1st (lines 221--234). Every definition, display and label of those lines is retained unchanged. Omitted from this package and not redistributed: 1--220, 235--271, 285--285, 353--978. Nothing inside a retained excerpt is omitted or altered: every retained body is delivered exactly as the article's lines are written, so no omission receipt of any kind accompanies this package. Citations: the retained excerpts carry no citation command at all, so no bibliography entry of the article is transcribed and no reference is redistributed. Reference adaptation: none was needed and none was made; every reference inside the delivered unit resolves inside it, so no unresolved reference or citation is produced. Printed slips kept literal: no printed word, symbol, hypothesis or number of the counted statement or of its proof was altered. Layout and class: the article's own class is \texttt{amsart} with packages including amsrefs; the wrapper is the lane's pinned \texttt{amsart} editorial wrapper at 10pt on A4, which restates the theorem-like environments the retained text needs and adds the article's own macro definitions that the retained text needs (\texttt{\textbackslash calU}, \texttt{\textbackslash calZ}, \texttt{\textbackslash cl}, \texttt{\textbackslash clbeta}, \texttt{\textbackslash preim}), with extra wrapper packages (bm, bbm, dsfont, xspace, cancel, mathtools). The delivered numbering is the unit's own. Assets: no retained span contains a figure environment, a graphics-inclusion command, a PDF-page inclusion, a table or a TikZ picture; no image file is copied into this package, the package declares no asset dependency and no image file is redistributed with it. Licence: version arXiv:2404.19703v3 of this article is distributed under the Creative Commons Attribution 4.0 International licence, as recorded in the arXiv licence metadata for this exact version and shown on the abstract page https://arxiv.org/abs/2404.19703v3. A full rights notice is delivered at licenses/arxiv-2404.19703-CC-BY-4.0.txt. Characters: every retained excerpt is pure ASCII, so the delivered engine, XeLaTeX with its default Latin Modern text font, reports no missing character.
\begin{lemma}\label{lemma.1st}
Let $A$ be a closed subset of a space $X$.
Then the following three conditions are equivalent.
\begin{enumerate}
\item $A$ is $C$-embedded in~$X$.
\item $A$ is $z$-embedded in~$X$ and for every zero-set $Z$ of~$\clbeta A$ 
      that is disjoint from~$A$ there is a zero-set~$Z^+$ of~$\beta X$ that 
      is disjoint from~$X$ and such that $Z=Z^+\cap\clbeta A$.
\item $A$ is $z$-embedded in~$X$ and for every zero-set $Z$ of~$\clbeta A$ 
      that is disjoint from~$A$ there is a countable family~$\calZ$ of 
      zero-sets of~$\beta X$ such that 
      $Z\subseteq\bigcup\calZ\subseteq\beta X\setminus X$.
\end{enumerate}
\end{lemma}

\begin{theorem}\label{thm.first}
Let $A$ be a closed subset of a space $X$.
Then the following three conditions are equivalent.
\begin{enumerate}
\item $A$ is Lindel\"of and $C$-embedded in~$X$.
\item For every compact subset $K$ of $\clbeta A\setminus A$ there is a 
      zero-set~$Z$ of~$\beta X$ such that 
      $K\subseteq Z\subseteq\beta X\setminus X$.
\item For every compact subset $K$ of $\clbeta A\setminus A$ there is a 
      countable family $\calZ$ of zero-sets of~$\beta X$ such that 
      $K\subseteq\bigcup\calZ\subseteq\beta X\setminus X$.
\end{enumerate}
\end{theorem}

\begin{proof}
To prove that (1) implies (2) we let $K$ be a compact subset 
of $\clbeta A\setminus A$ and find a zero-set $Z$ of $\beta X$ such 
that $K\subseteq Z\subseteq\beta X\setminus X$.

To begin we choose for every $a\in A$ a continuous function 
$f_a:\beta X\to[0,1]$ such that $f_a(a)=1$ and $f_a(x)=0$ if $x\in K$.
For each $a$ we let $U_a=f\preim\bigl[(\frac12,1]\bigr]$.
There is a countable subset $\{a_n:n\in\omega\}$ of~$A$ such 
that $A\subseteq\bigcup_{n\in\omega}U_{a_n}$.

Let $g=\sum_{n\in\omega}2^{-n}f_{a_n}$.
Then $g$~is continuous, $g(a)>0$ when $a\in A$, and $g(x)=0$ when $x\in K$.
We would like to let $Z=g\preim(0)$, but that set may intersect~$X$.
However, $S=\{x\in X:g(x)=0\}$ is a zero-set of~$X$ that is disjoint from~$A$.
As $A$~is $C$-embedded in~$X$ the sets $S$ and $A$ are completely separated.
Let $f:X\to[0,1]$ be continuous such that $f(x)=1$ if $x\in S$ and $f(a)=0$
if $a\in A$.
Note that $\beta f$ vanishes on $\clbeta A$ and in particular on~$K$.

Now let $h=g+\beta f$.
Then $h(x)\ge g(x)>0$ when $x\in X\setminus S$ and
$h(x)\ge1$ when $x\in S$.
Also, $h(x)=0$ when $x\in K$.
It follows that $h\preim(0)$ is the zero-set of~$\beta X$ that we seek.

\medskip
Clearly (2) implies (3).

\medskip
We finish by proving that (3) implies (1).   
To begin: the present condition~(3) is stronger than the second part of~(3)
in Lemma~\ref{lemma.1st}.
We need to show that $A$~is Lindel\"of and $z$-embedded in~$X$.
It will actually be simpler to show that $A$~is $C^*$-embedded in~$\clbeta A$.

\smallskip
That $A$~is Lindel\"of is proved as follows.
Let $\calU$ be a collection of open subsets of~$\beta X$ that covers~$A$.
Let $K=\clbeta A\setminus\bigcup\calU$ and let $\calZ$ be a countable
family of zero-sets of~$\beta X$ as in the assumption.

As in the proof of Lemma~\ref{lemma.1st} the complement~$L$ of
$\bigcup\calZ$ in~$\beta X$ is Lindel\"of and it contains~$X$.
Then $L\cap\clbeta A$ is Lindel\"of as well and, moreover,
contained in~$\bigcup\calU$.
But $A\subseteq L\cap\clbeta A$, so $A$~is covered by a countable
subfamily of~$\calU$.

\smallskip
To see that $A$~is $C^*$-embedded in~$\clbeta A$, and hence in~$\beta X$,
we show that if $E$ and $F$ are disjoint closed subsets of~$A$ then their 
closures in~$\clbeta A$ are disjoint; this shows that $\clbeta A$ actually
\emph{is} the \v{C}ech-Stone compactification of the normal space~$A$.

\smallskip
Let $E$ and~$F$ be as above and let $K=\clbeta E\cap\clbeta F$.
Then $K\subseteq\clbeta A\setminus A$ and hence there is a
countable family~$\calZ$ of zero-sets of~$\beta X$ as in our assumption.

We have just established that $L=\beta X\setminus\bigcup\calZ$ 
is Lindel\"of, hence $L$~is normal as well.
Also $X\subseteq L\subseteq\beta X$, and so $\beta L=\beta X$.

In addition we have $\cl_LE\cap\cl_LF=\emptyset$ and hence 
$\clbeta E\cap\clbeta F=\emptyset$ (so in hindsight $K=\emptyset$).
\end{proof}

\end{document}
